\chapter{Modèles}
\label{chap:modeles}

\section{Implémentation des modèles}
\label{sec:modeles}

Dans cette section, nous détaillons l'implémentation des différents modèles. 
Plusieurs approches seront comparées, en cohérence avec la littérature recensée, une 
approche statistiques (régression logistique, modèle de poisson), ensemblistes (Random Forest, Gradient 
Boosting). 
% Enfin, une approche par réseau de neurones, comme explorés par Carloni 
% \textit{et al.}\cite{carloni}.


Les performances seront évaluées sur un ensemble de test chronologique 
(saisons les plus récentes) afin d'éviter tout biais temporel, en utilisant 
des métriques adaptées : F1-score \ref{sec:f1}, Brier-Score \ref{sec:BS}, log-loss 
\ref{sec:log-loss} et ROI \ref{sec:roi}, que nous détaillerons dans \ref{chap:métriques}.

\subsection{Baseline}

En observant notre jeu de données, il est aisé de constater que prédire uniquement la 
victoire à domicile permet d'obtenir 45\% de bonne prédiction. Notre objectif est donc 
de dépasser ce seuil. 
 

\subsection{Régression Logistique}

La Régression Logistique constitue notre premier modèle de référence parmi les 
approches statistiques linéaires. Bien que simple, elle est particulièrement 
adaptée à la classification binaire (victoire à domicile vs. non-victoire). 
L'un de ses principaux atouts réside dans sa capacité à fournir des probabilités 
directement exploitables facilitant le calcul du ROI \ref{sec:roi}. 

Afin de maximiser ses performances, nous comparons une version dite de Base avec une 
version optimisée par Random Search\footnote{Le Random Search sélectionne au hasard des 
combinaisons d'hyperparamètres dans des plages prédéfinies afin de trouver la configuration 
la plus performante.}. 


\subsection{Random Forest}

Le Random Forest introduit une approche non linéaire basée sur le principe du Bagging 
combinant un grand nombre d'arbres de décision 
indépendants entraînés sur des sous-échantillons de données (bootstrap).

Dans notre pipeline, le Random Forest utilise des poids de classes équilibrés 
(class weight) pour compenser le déséquilibre initial du dataset (la victoire 
à domicile représentant 45\,\% des observations).
Nous comparons une version de base et une version 
optimisée par Random Search, portant sur le nombre d'arbres et la profondeur maximale. 

\subsection{XGBoost}

L'algorithme XGBoost (Extreme Gradient Boosting) représente l'état de l'art des 
méthodes ensemblistes par Boosting. Contrairement au Random Forest, les arbres 
de décision sont construits de manière séquentielle : chaque nouvel arbre est 
entraîné pour corriger les erreurs de prédiction (les résidus) commises par les 
précédents.

Nous implémentons une version de base et une version 
optimisée par Random Search. Pour cette dernière, l'optimisation des hyperparamètres 
(profondeur des arbres, taux d'apprentissage \textit{learning rate}, et paramètres 
de régularisation $L_1$/$L_2$) est couplée à une validation croisée temporelle, 
garantissant la robustesse des prédictions sur les saisons futures.

\subsection{SVM (Support Vector Machine)}

Le SVM à noyau de fonction de base radiale (RBF) adopte une approche géométrique. 
Il projette notre espace de caractéristiques de forme et de classement dans une 
dimension supérieure afin de trouver l'hyperplan de séparation optimal à marge 
maximale entre les matchs gagnés et perdus.

Le noyau RBF est très efficace lorsque la frontière de décision est fortement 
non linéaire. Cependant, le SVM ne fournit pas nativement de probabilités, mais 
des scores de distance à l'hyperplan. Pour intégrer ce modèle nous activons la méthode 
de calibration de Platt (\texttt{probability=True}), 
convertissant ces scores en probabilités. Cela nous permet d'évaluer la fiabilité 
de ses prédictions au même titre que les autres algorithmes.

\subsection{LightGBM}

Enfin, nous avons également utilisé LightGBM \cite{chenlgbm} 
(Light Gradient Boosting Machine), c'est une déclinaison optimisée 
du Gradient Boosting développée par Microsoft. Sa spécificité 
repose sur une croissance des arbres par feuille 
(leaf-wise) plutôt que par niveau (level-wise), ce qui lui confère une 
vitesse de convergence supérieure et une excellente gestion des interactions 
de variables à grande échelle.



% \subsection{Réseau de neurones}
